\documentclass[11pt]{article}

\usepackage[margin=1in]{geometry}
\usepackage{amsmath,amssymb,amsthm}
\usepackage{graphicx}
\usepackage[hidelinks]{hyperref}

\newtheorem{theorem}{Theorem}
\newtheorem{lemma}{Lemma}
\newtheorem{corollary}{Corollary}
\newtheorem{remark}{Remark}

\newcommand{\eps}{\varepsilon}
\newcommand{\dist}{\operatorname{dist}}
\newcommand{\diam}{\operatorname{diam}}

\title{The ball dentable property for Bochner \(L^p\)-spaces}
\author{Agent lane 04}
\date{2026-06-28}

\begin{document}
\maketitle

\begin{abstract}
Basu and Seal prove that \(L^p(\mu,X)\) having the ball dentable property implies that \(X\) has the ball dentable property, and ask whether the converse holds.  We prove the converse for \(1<p<\infty\) and finite nonzero \(\mu\).  Thus, in this range, \(L^p(\mu,X)\) has BDP if and only if \(X\) has BDP.  Combined with the scalar \(L^1[0,1]\) endpoint counterexample already recorded in the run, this gives the BDP exponent split for the Bochner converse: false at \(p=1\), true for \(1<p<\infty\).
\end{abstract}

\section{Source Question}

The source paper introduces the ball dentable property (BDP): a Banach space \(X\) has BDP if, for every \(\eps>0\), the unit ball \(B_X\) has a slice of diameter less than \(\eps\).

In Section 3, the source proposition labelled \texttt{bochner bdp} proves:
\[
  L^p(\mu,X)\text{ has BDP}\quad\Longrightarrow\quad X\text{ has BDP}
\]
for \(1\le p<\infty\) and finite \(\mu\).  Immediately after the proof, the authors write that they do not know whether the converse is true in general.  The relevant source passage is:

\begin{center}
\includegraphics[width=.92\linewidth]{figures/open_problem_crop.png}
\end{center}

The \(p=1\) converse is already false, even for \(X=\mathbb R\) and Lebesgue measure, because every slice of the unit ball of non-atomic \(L^1\) has diameter two.  The present note proves the positive \(1<p<\infty\) half for BDP.

\section{Idea of the Proof}

Let \(S(B_X,x^*,\eta)\) be a small slice of \(B_X\).  We lift it to a slice of \(B_{L^p(\mu,X)}\) by the constant dual function
\[
  F(f)=\int_\Omega x^*(f(t))\,d\mu(t)
\]
after normalizing \(\mu\) to be a probability measure.

The key point is specific to \(1<p<\infty\).  In scalar \(L^p\), the functional \(h\mapsto\int h\) strongly exposes the constant function \(1\) on the positive part of the unit ball.  Thus if \(F(f)\) is close to \(1\), then \(\|f(\cdot)\|\) is close to \(1\) in \(L^p\), and the nonnegative defect
\[
  \|f(t)\|-x^*(f(t))
\]
has small integral.  Except on a set of small \(L^p\)-mass, the normalized vector \(f(t)/\|f(t)\|\) must lie in the original small slice of \(B_X\).  Two functions in the lifted slice are therefore pointwise close on most of \(\Omega\), and the exceptional parts have small \(L^p\)-norm.

\section{Proof}

We use the notation
\[
  S(B_X,x^*,\eta)=\{x\in B_X: x^*(x)>1-\eta\},
\]
where \(x^*\in S_{X^*}\) and \(\eta>0\).

\begin{lemma}[Lifting a small BDP slice]\label{lem:lift}
Let \((\Omega,\Sigma,\mu)\) be a probability space, let \(1<p<\infty\), and let \(x^*\in S_{X^*}\).  Suppose that
\[
  S=S(B_X,x^*,\eta)
\]
has diameter \(d\).  For every \(\gamma>0\) there is \(\alpha>0\) such that the slice
\[
  \mathcal S_\alpha
  =
  \left\{f\in B_{L^p(\mu,X)}:
  \int_\Omega x^*(f(t))\,d\mu(t)>1-\alpha\right\}
\]
has diameter at most \(d+\gamma\).
\end{lemma}

\begin{proof}
It is enough to prove the following sequential assertion: if \(\alpha_n\downarrow0\) and \(f_n,g_n\in B_{L^p(\mu,X)}\) satisfy
\[
  \int x^*(f_n)\,d\mu>1-\alpha_n,
  \qquad
  \int x^*(g_n)\,d\mu>1-\alpha_n,
\]
then
\[
  \limsup_{n\to\infty}\|f_n-g_n\|_{L^p(\mu,X)}\le d .
\]
Indeed, failure of the lemma for some \(\gamma>0\) would give such a sequence with all distances \(>d+\gamma\).

Put \(r_n(t)=\|f_n(t)\|\) and \(s_n(t)=\|g_n(t)\|\).  Since
\[
  \int x^*(f_n)\,d\mu
  \le \int r_n\,d\mu
  \le \|r_n\|_{L^p}
  \le 1,
\]
we have \(\int r_n\,d\mu\to1\) and \(\|r_n\|_{L^p}\to1\).  The scalar space \(L^p(\mu)\) is uniformly convex, so the constant function \(1\) is strongly exposed in \(B_{L^p(\mu)}\) by the functional \(h\mapsto\int h\,d\mu\).  Hence
\[
  r_n\to 1\quad\text{in }L^p(\mu).
\]
The same argument gives \(s_n\to1\) in \(L^p(\mu)\).  Moreover,
\[
  \int_\Omega (r_n(t)-x^*(f_n(t)))\,d\mu(t)\to0,
  \qquad
  \int_\Omega (s_n(t)-x^*(g_n(t)))\,d\mu(t)\to0.
\]

Define the bad sets
\[
  A_n=\{t:r_n(t)=0\}\cup\left\{t:r_n(t)>0\text{ and }
  x^*\!\left(\frac{f_n(t)}{r_n(t)}\right)\le1-\eta\right\},
\]
and similarly
\[
  B_n=\{t:s_n(t)=0\}\cup\left\{t:s_n(t)>0\text{ and }
  x^*\!\left(\frac{g_n(t)}{s_n(t)}\right)\le1-\eta\right\}.
\]
On the part of \(A_n\) where \(r_n>0\) we have
\[
  r_n(t)-x^*(f_n(t))
  =
  r_n(t)\left(1-x^*\!\left(\frac{f_n(t)}{r_n(t)}\right)\right)
  \ge \eta r_n(t),
\]
so \(\int_{A_n} r_n\,d\mu\to0\), while the zero part of \(A_n\) has measure at most \(\|r_n-1\|_p^p\).  Since \(r_n\to1\) in \(L^p\), this implies \(\mu(A_n)\to0\) and then
\[
  \int_{A_n} r_n^p\,d\mu\to0.
\]
For example, \(A_n\cap\{r_n>1/2\}\) has measure at most \(2\int_{A_n}r_n\,d\mu\), while \(A_n\cap\{r_n\le1/2\}\) has measure controlled by \(2^p\|r_n-1\|_p^p\); after this, \(r_n^p\le 2^{p-1}(|r_n-1|^p+1)\) gives the displayed convergence.  The same reasoning gives
\[
  \mu(B_n)\to0,\qquad \int_{B_n}s_n^p\,d\mu\to0.
\]

Let \(C_n=\Omega\setminus(A_n\cup B_n)\).  On \(C_n\), both \(r_n\) and \(s_n\) are positive, and the normalized vectors
\[
  u_n(t)=\frac{f_n(t)}{r_n(t)},\qquad
  v_n(t)=\frac{g_n(t)}{s_n(t)}
\]
belong to the slice \(S\).  Hence \(\|u_n(t)-v_n(t)\|\le d\) on \(C_n\).  Therefore
\[
\begin{aligned}
  \|(f_n-g_n)\mathbf 1_{C_n}\|_p
  &\le
  \|(r_n-s_n)\mathbf 1_{C_n}\|_p
  + d\,\|s_n\mathbf 1_{C_n}\|_p  \\
  &\le
  \|r_n-1\|_p+\|s_n-1\|_p+d.
\end{aligned}
\]
The right-hand side has limsup at most \(d\).

It remains to check that the complement contributes nothing in the limit.  Since \(r_n\to1\) and \(s_n\to1\) in \(L^p\), and since \(\mu(A_n\cup B_n)\to0\), we have
\[
  \int_{A_n\cup B_n} r_n^p\,d\mu\to0,\qquad
  \int_{A_n\cup B_n} s_n^p\,d\mu\to0.
\]
Thus
\[
  \|(f_n-g_n)\mathbf 1_{\Omega\setminus C_n}\|_p
  \le
  \|r_n\mathbf 1_{A_n\cup B_n}\|_p
  +
  \|s_n\mathbf 1_{A_n\cup B_n}\|_p
  \to0.
\]
Combining the two estimates proves
\[
  \limsup_n \|f_n-g_n\|_p\le d,
\]
as required.
\end{proof}

\begin{theorem}\label{thm:bdp-positive}
Let \((\Omega,\Sigma,\mu)\) be a finite nonzero measure space, let \(1<p<\infty\), and let \(X\) be a Banach space.  If \(X\) has BDP, then \(L^p(\mu,X)\) has BDP.
\end{theorem}

\begin{proof}
Since replacing \(\mu\) by the probability measure \(\mu/\mu(\Omega)\) only rescales the \(L^p\)-norm by a constant, we may assume that \(\mu\) is a probability measure.

Fix \(\eps>0\).  Since \(X\) has BDP, choose \(x^*\in S_{X^*}\) and \(\eta>0\) such that the slice
\[
  S(B_X,x^*,\eta)
\]
has diameter less than \(\eps/2\).  Apply Lemma \ref{lem:lift} with \(\gamma=\eps/2\).  The resulting slice \(\mathcal S_\alpha\) of \(B_{L^p(\mu,X)}\) has diameter less than \(\eps\).  Therefore \(L^p(\mu,X)\) has BDP.
\end{proof}

\begin{corollary}\label{cor:equivalence}
Let \((\Omega,\Sigma,\mu)\) be a finite nonzero measure space and \(1<p<\infty\).  Then
\[
  L^p(\mu,X)\text{ has BDP}
  \quad\Longleftrightarrow\quad
  X\text{ has BDP}.
\]
\end{corollary}

\begin{proof}
The forward implication is the proposition labelled \texttt{bochner bdp} in Basu--Seal.  The reverse implication is Theorem \ref{thm:bdp-positive}.
\end{proof}

\section{Endpoint and Scope}

The theorem is sharp in the exponent.  At \(p=1\), the converse fails over the Lebesgue interval with \(X=\mathbb R\): the space \(X\) has BDP, while every slice of the unit ball of scalar \(L^1[0,1]\) has diameter two.  This endpoint counterexample is recorded separately in the run packet
\[
  \texttt{2011.14591\_l1\_scalar\_bdp\_converse\_counterexample}.
\]

This packet does not address the analogous BHP converse.  The proof above uses a slice and the strong exposure of the scalar \(L^p\) unit ball; it does not convert an arbitrary small relatively weakly open subset of \(B_X\) into a small relatively weakly open subset of \(B_{L^p(\mu,X)}\).

\section{Novelty Check}

The cheap run indexes were searched for arXiv:2108.02908, arXiv:2011.14591, BDP, BHP, Lebesgue-Bochner, Ball Dentable Property, and the converse language from the source.  The existing \(p=1\) endpoint packet was found and is cited above; it explicitly left the \(1<p<\infty\) positive stability direction outside its scope.

Web searches for exact phrases including ``Ball Dentable Property \(L^p\) Bochner'', ``BDP Lebesgue-Bochner \(p>1\)'', and the source title returned the source paper but did not locate a separate explicit answer to the \(1<p<\infty\) BDP converse.

\section{References}

\begin{itemize}
\item S. Basu and S. Seal, \emph{Small Combination of Slices, Dentability and Stability Results Of Small Diameter Properties In Banach Spaces}, arXiv:2108.02908.
\item S. Basu and S. Seal, \emph{Stability results of small diameter properties in Banach spaces}, arXiv:2011.14591.  This is the related arXiv record used by the existing endpoint packet.
\end{itemize}

\end{document}
